\scanpage{65}
\mspage{14}

\deleted{\unsure{d|b}}\unsure{b|d}$''$) Soit $W_\omega\deleted{\unsure{\in|=}}\subset W\cap W_\infty$ la plus petite partie de $\mathrm{Fl}(\mathrm{Cat})$ satisfaisant les conditions $\textcircled{1}\textcircled{2}\textcircled{3}{}^*$ (axiome de saturation, de localisation, de l'objet final). Soit $f:X\to Y$ une flèche de $\mathrm{Cat}$ qui soit $W_\omega$-parfaite. Si $f\in W$, alors les fibres de $f$ sont $W_\omega$-asphériques.

% slanted
\margin{Je vais même prendre $\text{\textcircled{\scriptsize 2 bis}}$ au lieu de $\textcircled{2}$, pour trouver $W_\omega$ le plus petit possible !}

Entre ces propriétés de $W$, on a les \unsure{relations|implications} tautologiques (sauf $a''\Rightarrow b)$, qui est \unsure{exact.} la prop. 4)

\deleted{\diagram{\begin{tikzcd}[ampersand replacement=\&,row sep=small,column sep=tiny]
\unsure{a_1|a}\arrow[dr,Rightarrow] \& \unsure{a|a_1}\arrow[d,Rightarrow]\arrow[r,Leftrightarrow] \& \unsure{a_2|a}\arrow[dl,Rightarrow]\\
\& b\arrow[d,Rightarrow]\\
\& c\arrow[d,Rightarrow]\\
\& d\arrow[d,Rightarrow]\\
\& d'\arrow[d,Rightarrow]\\
\& d''.
\end{tikzcd}}}

\diagram{\begin{tikzcd}[ampersand replacement=\&,row sep=small]
a=(a_1\Leftrightarrow a_2)\arrow[d,Rightarrow]\\
a'\arrow[d,Rightarrow]\\
a''\arrow[d,Rightarrow]\\
b\arrow[d,Rightarrow]\\
b'\arrow[d,Rightarrow]\\
b''
\end{tikzcd}}

\margin{$\left.\vphantom{\begin{matrix}a\\a'\\a''\end{matrix}}\right\}$ En fait, toutes ces conditions sont équivalentes.}

% slanted
\margin{Je reviens sur les diverses variantes de $W(4)$ plus bas, p. \unsure{53-60|53-66|53-56}}

Eh bien, dans le lemme page 12, il suffisait de supposer que $W$ satisfait à la version la plus faible $d'')$ de l'axiome de fibration (sept versions !). En effet, les fibrations $\bar X\unsure{\mathrel{-}|\to}Y$ et $E\to X$ utilisées dans la démonstration du lemme sont $W_\omega$-parfaites, et il suffisait de savoir que si elles sont $\in W$, alors les fibres sont $W$-asphériques.

Mais je ne pense pas que même la version la plus forte $4\,\mathrm a)$ de \deleted{\unsure{fibra}} l'axiome de fibration, jointe aux trois axiomes précédents\inserted{et à 5 (connexité)}, implique $W=W_\infty$. Il faudrait tout au moins y joindre l'axiome.

% slanted
\margin{$W(6)$ Axiome des limites filtrantes}

$\textcircled{6}$ Soit $X=\underset{i}{\underline{\lim}}X_i$ une limite inductive filtrante dans $\mathrm{Cat}$. Si les $X_i$ sont $W$-asphériques, $X$ aussi.

Ceci devrait être équivalent à la propriété, en apparence plus forte, en termes de \deleted{\unsure{toutes}} flèches $\in W$ :
\uline{Proposition \unsure{5}} Supposons qu'on ait$^{\unsure{1|\wedge}}$ $\textcircled{6}$, (en plus de $\textcircled{1}$ à $\textcircled{4}{}^*$),

% slanted
\margin{$\textcircled{4}$ sous la forme la plus faible \deleted{\unsure{$4\,\mathrm b$}} $\text{\textcircled{\scriptsize $4\,\mathrm d''$}}$}

% archive sheet number: 65
